% Module 3: Qubit States, Gates and Circuits
% Estimated slides: 45 | Duration: ~3 hours
% Key topics: Quantum state notation, single-qubit gates (X/Y/Z/H/S/T/Rx/Ry/Rz),
%             circuit diagrams, multi-qubit gates (CNOT/Toffoli/SWAP), Bell states,
%             circuit depth/width, universality
% Labs: Single-qubit circuits in Qiskit; Bell state preparation

%%%%%%%%%%%%%%%%%%%%%%%%%%%%%%%%%%%%%%%%%%%%%%%%%%%%%%%%%%
\begin{frame}[fragile]\frametitle{}
\begin{center}
{\Large Module 3: Qubit States, Gates and Circuits}
\end{center}
\end{frame}

%%%%%%%%%%%%%%%%%%%%%%%%%%%%%%%%%%%%%%%%%%%%%%%%%%%%%%%%%%%
\begin{frame}[fragile]\frametitle{Quantum State Notation (Minimal)}
\begin{itemize}
\item We write qubit states using Dirac notation: $|0\rangle$ and $|1\rangle$
\item Read $|0\rangle$ as ``ket zero'' and $|1\rangle$ as ``ket one''
\item A general qubit state: $|\psi\rangle = \alpha|0\rangle + \beta|1\rangle$
\item $|\alpha|^2$ = probability of measuring 0; $|\beta|^2$ = probability of measuring 1
\item $|+\rangle = (|0\rangle + |1\rangle)/\sqrt{2}$: equal superposition
\item $|-\rangle = (|0\rangle - |1\rangle)/\sqrt{2}$: equal superposition but with phase flip
\item This notation is convenient --- we will use it throughout
\end{itemize}
% Suggested diagram: Bloch sphere with key states labeled in ket notation
\end{frame}

%%%%%%%%%%%%%%%%%%%%%%%%%%%%%%%%%%%%%%%%%%%%%%%%%%%%%%%%%%%
\begin{frame}[fragile]\frametitle{Quantum Gates: The Basics}
\begin{itemize}
\item A quantum gate transforms one qubit state into another
\item Represented mathematically as a matrix (unitary matrix)
\item Key property: quantum gates are \textit{reversible} --- always invertible
\item Circuit symbols: boxes or standard symbols on wires representing qubit evolution
\item Reading circuit diagrams: time flows left to right, qubits are horizontal wires
\item Gates are applied top-to-bottom on each wire
\end{itemize}
% Suggested diagram: Simple 2-qubit circuit with gate boxes labeled
\end{frame}

%%%%%%%%%%%%%%%%%%%%%%%%%%%%%%%%%%%%%%%%%%%%%%%%%%%%%%%%%%%
\begin{frame}[fragile]\frametitle{The Pauli-X Gate: Quantum NOT}
\begin{itemize}
\item X gate = quantum NOT gate; flips $|0\rangle \to |1\rangle$ and $|1\rangle \to |0\rangle$
\item On the Bloch sphere: rotation by 180 degrees around the X axis
\item Matrix: $\begin{pmatrix}0 & 1 \\ 1 & 0\end{pmatrix}$
\item Circuit symbol: a box with ``X'' or a circled plus $\oplus$
\item Qiskit: \texttt{qc.x(qubit)}
\item Example: $X|0\rangle = |1\rangle$, $X|+\rangle = |+\rangle$ (superposition is unchanged)
\end{itemize}
% Suggested diagram: Bloch sphere showing X rotation; circuit symbol
\end{frame}

%%%%%%%%%%%%%%%%%%%%%%%%%%%%%%%%%%%%%%%%%%%%%%%%%%%%%%%%%%%
\begin{frame}[fragile]\frametitle{The Pauli-Y and Pauli-Z Gates}
\begin{itemize}
\item \textbf{Y gate}: rotation by 180 degrees around Y axis; $|0\rangle \to i|1\rangle$, $|1\rangle \to -i|0\rangle$
\item Y = combination of X and Z with a phase; less common in circuits
\item \textbf{Z gate}: rotation around Z axis; $|0\rangle \to |0\rangle$, $|1\rangle \to -|1\rangle$
\item Z flips the phase of $|1\rangle$ but leaves $|0\rangle$ unchanged
\item Z leaves the measurement probabilities unchanged but affects interference
\item Qiskit: \texttt{qc.y(qubit)}, \texttt{qc.z(qubit)}
\end{itemize}
% Suggested diagram: Bloch sphere showing Y and Z rotations; circuit symbols
\end{frame}

%%%%%%%%%%%%%%%%%%%%%%%%%%%%%%%%%%%%%%%%%%%%%%%%%%%%%%%%%%%
\begin{frame}[fragile]\frametitle{The Hadamard Gate: Superposition Creator}
\begin{itemize}
\item H gate transforms $|0\rangle \to |+\rangle$ and $|1\rangle \to |-\rangle$
\item Starting from $|0\rangle$, H creates equal superposition: 50\% chance of 0 or 1
\item The most commonly used gate in quantum algorithms
\item Applying H twice returns to the original state: $H \cdot H = I$
\item On Bloch sphere: rotation by 180 degrees around the axis halfway between X and Z
\item Qiskit: \texttt{qc.h(qubit)}
\end{itemize}
% Suggested diagram: Bloch sphere before and after H gate; circuit symbol H
\end{frame}

%%%%%%%%%%%%%%%%%%%%%%%%%%%%%%%%%%%%%%%%%%%%%%%%%%%%%%%%%%%
\begin{frame}[fragile]\frametitle{Phase Gates: S and T}
\begin{itemize}
\item Phase gates add a phase angle to the $|1\rangle$ component without changing probabilities
\item \textbf{S gate}: adds a $90°$ phase to $|1\rangle$; $|1\rangle \to i|1\rangle$
\item \textbf{T gate}: adds a $45°$ phase to $|1\rangle$; $|1\rangle \to e^{i\pi/4}|1\rangle$
\item Phase gates look invisible if you only measure in the Z basis
\item They become important when combined with H gates (interference)
\item Qiskit: \texttt{qc.s(qubit)}, \texttt{qc.t(qubit)}
\end{itemize}
% Suggested diagram: Bloch sphere rotation around Z axis for S and T
\end{frame}

%%%%%%%%%%%%%%%%%%%%%%%%%%%%%%%%%%%%%%%%%%%%%%%%%%%%%%%%%%%
\begin{frame}[fragile]\frametitle{Rotation Gates: Rx, Ry, Rz}
\begin{itemize}
\item Generalize Pauli rotations to arbitrary angles
\item $R_x(\theta)$: rotate by angle $\theta$ around X axis on Bloch sphere
\item $R_y(\theta)$: rotate by angle $\theta$ around Y axis
\item $R_z(\theta)$: rotate by angle $\theta$ around Z axis (generalized phase)
\item Used in variational/parameterized circuits (VQE, QAOA)
\item Qiskit: \texttt{qc.rx(theta, qubit)}, \texttt{qc.ry(theta, qubit)}, \texttt{qc.rz(theta, qubit)}
\item Special cases: $R_x(\pi) = X$, $R_z(\pi) = Z$ (up to global phase)
\end{itemize}
% Suggested diagram: Parameterized rotation on Bloch sphere
\end{frame}

%%%%%%%%%%%%%%%%%%%%%%%%%%%%%%%%%%%%%%%%%%%%%%%%%%%%%%%%%%%
\begin{frame}[fragile]\frametitle{Single-Qubit Gate Summary Table}
\begin{itemize}
\item Quick reference for all single-qubit gates
\end{itemize}
\begin{center}
% Suggested diagram: Table with columns: Gate | Symbol | Action | Bloch rotation | Use case
% X | X | Flips 0 to 1 | 180 deg X-axis | NOT operation
% Z | Z | Phase flip of |1> | 180 deg Z-axis | Phase kick
% H | H | Superposition | 180 deg XZ diagonal | Start algorithms
% S | S | 90 deg phase | 90 deg Z-axis | QFT
% T | T | 45 deg phase | 45 deg Z-axis | Universal gate set
\end{center}
\end{frame}

%%%%%%%%%%%%%%%%%%%%%%%%%%%%%%%%%%%%%%%%%%%%%%%%%%%%%%%%%%%
\begin{frame}[fragile]\frametitle{Reading Quantum Circuit Diagrams}
\begin{itemize}
\item Each horizontal line = one qubit (initialized to $|0\rangle$ unless noted)
\item Each box on a wire = a gate applied to that qubit
\item Time flows left to right
\item A double horizontal line = classical bit wire (after measurement)
\item Measurement symbol: meter icon, connects to classical wire
\item Vertical line connecting wires = two-qubit gate
\item Circuit width = number of qubits; circuit depth = longest gate sequence
\end{itemize}
% Suggested diagram: Annotated circuit diagram with labels pointing to each element
\end{frame}

%%%%%%%%%%%%%%%%%%%%%%%%%%%%%%%%%%%%%%%%%%%%%%%%%%%%%%%%%%%
\begin{frame}[fragile]\frametitle{Lab: Single-Qubit Circuits in Qiskit}
\begin{itemize}
\item Build circuits applying different single-qubit gates
\item Visualize with \texttt{qc.draw(output='mpl')}
\item Measure and plot histograms
\item Task 1: $|0\rangle \to X \to$ measure. Expected: always 1
\item Task 2: $|0\rangle \to H \to$ measure. Expected: 50/50
\item Task 3: $|0\rangle \to H \to Z \to H \to$ measure. Expected: always 1 (why?)
\end{itemize}
\begin{verbatim}
qc = QuantumCircuit(1, 1)
qc.h(0); qc.z(0); qc.h(0)
qc.measure(0, 0)
qc.draw('mpl')
\end{verbatim}
\end{frame}

%%%%%%%%%%%%%%%%%%%%%%%%%%%%%%%%%%%%%%%%%%%%%%%%%%%%%%%%%%%
\begin{frame}[fragile]\frametitle{Multi-Qubit Gates: The CNOT Gate}
\begin{itemize}
\item CNOT = Controlled-NOT: flips the target qubit only if the control qubit is $|1\rangle$
\item $|00\rangle \to |00\rangle$, $|01\rangle \to |01\rangle$, $|10\rangle \to |11\rangle$, $|11\rangle \to |10\rangle$
\item The quantum equivalent of the classical XOR gate
\item Most important two-qubit gate: creates entanglement when combined with H
\item Circuit symbol: filled circle on control qubit, $\oplus$ on target qubit
\item Qiskit: \texttt{qc.cx(control, target)}
\end{itemize}
% Suggested diagram: CNOT gate truth table and circuit symbol with control and target
\end{frame}

%%%%%%%%%%%%%%%%%%%%%%%%%%%%%%%%%%%%%%%%%%%%%%%%%%%%%%%%%%%
\begin{frame}[fragile]\frametitle{Creating Entanglement: H then CNOT}
\begin{itemize}
\item Start with two qubits in $|00\rangle$
\item Apply H to qubit 0: qubit 0 becomes $|+\rangle = (|0\rangle + |1\rangle)/\sqrt{2}$
\item Combined state: $(|0\rangle + |1\rangle)/\sqrt{2} \otimes |0\rangle = (|00\rangle + |10\rangle)/\sqrt{2}$
\item Apply CNOT: $|10\rangle \to |11\rangle$, giving $(|00\rangle + |11\rangle)/\sqrt{2}$
\item This is the Bell state $|\Phi^+\rangle$ --- maximally entangled
\item Measure: always get 00 or 11, never 01 or 10
\end{itemize}
% Suggested diagram: Step-by-step state evolution for Bell state creation
\end{frame}

%%%%%%%%%%%%%%%%%%%%%%%%%%%%%%%%%%%%%%%%%%%%%%%%%%%%%%%%%%%
\begin{frame}[fragile]\frametitle{The Four Bell States}
\begin{itemize}
\item $|\Phi^+\rangle = (|00\rangle + |11\rangle)/\sqrt{2}$: qubit outcomes always equal
\item $|\Phi^-\rangle = (|00\rangle - |11\rangle)/\sqrt{2}$: equal outcomes, phase flip
\item $|\Psi^+\rangle = (|01\rangle + |10\rangle)/\sqrt{2}$: qubit outcomes always opposite
\item $|\Psi^-\rangle = (|01\rangle - |10\rangle)/\sqrt{2}$: opposite outcomes, phase flip
\item These 4 states form a basis for the 2-qubit space
\item Used in teleportation, superdense coding, QKD
\end{itemize}
% Suggested diagram: 2x2 grid of Bell states with circuits to prepare each
\end{frame}

%%%%%%%%%%%%%%%%%%%%%%%%%%%%%%%%%%%%%%%%%%%%%%%%%%%%%%%%%%%
\begin{frame}[fragile]\frametitle{The Toffoli Gate (CCNOT)}
\begin{itemize}
\item Toffoli = Doubly Controlled NOT: flips target only if BOTH control qubits are $|1\rangle$
\item Quantum version of the classical AND gate
\item Toffoli gate is universal for classical computation within quantum circuits
\item Circuit symbol: two control circles and a target $\oplus$
\item Qiskit: \texttt{qc.ccx(ctrl1, ctrl2, target)}
\item Used in Grover's oracle, arithmetic circuits, error correction
\end{itemize}
% Suggested diagram: Toffoli gate truth table and circuit symbol
\end{frame}

%%%%%%%%%%%%%%%%%%%%%%%%%%%%%%%%%%%%%%%%%%%%%%%%%%%%%%%%%%%
\begin{frame}[fragile]\frametitle{The SWAP Gate}
\begin{itemize}
\item SWAP exchanges the states of two qubits
\item $|01\rangle \to |10\rangle$ and $|10\rangle \to |01\rangle$
\item Used to move quantum information between qubits (routing)
\item On hardware with limited connectivity, SWAP gates are inserted during transpilation
\item Expensive: SWAP = 3 CNOT gates
\item Qiskit: \texttt{qc.swap(qubit1, qubit2)}
\end{itemize}
% Suggested diagram: SWAP gate circuit symbol; decomposition into 3 CNOTs
\end{frame}

%%%%%%%%%%%%%%%%%%%%%%%%%%%%%%%%%%%%%%%%%%%%%%%%%%%%%%%%%%%
\begin{frame}[fragile]\frametitle{Controlled Gates: General Pattern}
\begin{itemize}
\item Any single-qubit gate U can be made controlled: apply U to target only if control is $|1\rangle$
\item CX (CNOT), CZ, CH, CS, CT are all controlled versions of X, Z, H, S, T
\item Generalization: multi-controlled gates (ccx, cccx, ...)
\item Qiskit: \texttt{qc.cx(c,t)}, \texttt{qc.cz(c,t)}, \texttt{qc.ch(c,t)}, \texttt{qc.cp(theta, c, t)}
\item Controlled gates are the primary way to create entanglement
\end{itemize}
% Suggested diagram: General controlled-U gate symbol
\end{frame}

%%%%%%%%%%%%%%%%%%%%%%%%%%%%%%%%%%%%%%%%%%%%%%%%%%%%%%%%%%%
\begin{frame}[fragile]\frametitle{Circuit Depth and Width}
\begin{itemize}
\item \textbf{Width}: number of qubits = computational resource space
\item \textbf{Depth}: number of sequential gate layers = time resource
\item Gates that operate on different qubits can run in parallel (same layer)
\item NISQ constraint: depth limited by decoherence time / T2
\item Rule of thumb for current hardware: aim for depth < 100 gates
\item Transpilation may increase depth when converting to native gate set
\end{itemize}
% Suggested diagram: Circuit showing depth=3, width=3, with parallel gates in layers
\end{frame}

%%%%%%%%%%%%%%%%%%%%%%%%%%%%%%%%%%%%%%%%%%%%%%%%%%%%%%%%%%%
\begin{frame}[fragile]\frametitle{Universal Gate Sets}
\begin{itemize}
\item A gate set is universal if any unitary operation can be approximated with it
\item Classical analogy: NAND gate is universal for classical logic
\item Quantum universal sets: $\{H, T, CNOT\}$ is sufficient
\item In practice, hardware has its own native gate set (e.g., IBM: CX, $R_z$, $\sqrt{X}$)
\item Transpiler converts your circuit gates to the hardware's native set
\item More gates = more errors; decomposition adds overhead
\end{itemize}
% Suggested diagram: Gate set hierarchy and decomposition example
\end{frame}

%%%%%%%%%%%%%%%%%%%%%%%%%%%%%%%%%%%%%%%%%%%%%%%%%%%%%%%%%%%
\begin{frame}[fragile]\frametitle{Lab: Bell State Preparation in Qiskit}
\begin{itemize}
\item Build all four Bell states using QuantumCircuit
\item Verify each with statevector simulation and measurement histogram
\item Key observation: perfect anti-correlation or correlation
\item Extension: modify the Bell state creation to add a Z gate before CNOT and observe the difference
\end{itemize}
\begin{verbatim}
from qiskit.quantum_info import Statevector
qc = QuantumCircuit(2)
qc.h(0); qc.cx(0, 1)
sv = Statevector(qc)
print(sv)  # Should show (|00>+|11>)/sqrt(2)
\end{verbatim}
\end{frame}

%%%%%%%%%%%%%%%%%%%%%%%%%%%%%%%%%%%%%%%%%%%%%%%%%%%%%%%%%%%
\begin{frame}[fragile]\frametitle{Lab: Circuit Diagram Reading Exercise}
\begin{itemize}
\item Given circuit: $|0\rangle \otimes |0\rangle$, then H on q0, CX(q0, q1), H on q0, measure both
\item Predict the output state before measuring
\item Build and run in Qiskit to verify
\item Extension: add a Z gate between H and CX on q0 --- what changes?
\end{itemize}
% Suggested diagram: The circuit described above drawn as a diagram for students to analyze
\end{frame}

%%%%%%%%%%%%%%%%%%%%%%%%%%%%%%%%%%%%%%%%%%%%%%%%%%%%%%%%%%%
\begin{frame}[fragile]\frametitle{Exercise: Gate Identification}
\begin{itemize}
\item Match each operation to the correct gate:
\begin{enumerate}
  \item Flip $|0\rangle$ to $|1\rangle$ and $|1\rangle$ to $|0\rangle$: (X gate)
  \item Create equal superposition from $|0\rangle$: (H gate)
  \item Flip target qubit only if control is 1: (CNOT/CX gate)
  \item Add a 90-degree phase to $|1\rangle$: (S gate)
  \item Create entanglement between two qubits: (H + CNOT)
  \item Flip target only if two controls are both 1: (Toffoli gate)
\end{enumerate}
\end{itemize}
\end{frame}

%%%%%%%%%%%%%%%%%%%%%%%%%%%%%%%%%%%%%%%%%%%%%%%%%%%%%%%%%%%
\begin{frame}[fragile]\frametitle{Module 3 Summary}
\begin{itemize}
\item Quantum gates are unitary (reversible) transformations on qubit states
\item Key single-qubit gates: X (NOT), H (superposition), Z (phase), S, T, Rx, Ry, Rz
\item Key multi-qubit gates: CNOT (entanglement), Toffoli (AND), SWAP
\item Circuit diagrams: qubits as wires, gates as boxes, time flows left to right
\item Bell states: maximally entangled 2-qubit states, prepared with H+CNOT
\item Universal gate set: $\{H, T, CNOT\}$ can approximate any quantum operation
\item Circuit depth limited by decoherence --- keep circuits shallow on NISQ devices
\item Next: Qiskit programming lab (Workshop B, Module 11)
\end{itemize}
\end{frame}
